\chapter{Реализованные шаблоны prompts}

В этом приложении фиксируются основные шаблоны, уже реализованные в репозитории. Они включены сюда, поскольку задают операциональную форму prompts, используемых на протяжении пилотного эксперимента.

\section{Prompt для генерации референсов}

Prompt, используемый для reference retrieval и candidate generation:
\begin{quote}
    \ttfamily
    Write a joke using the following keyword(s): \{keywords\}
\end{quote}

Его главное преимущество --- согласованность. Одна и та же текстовая форма используется для embedding prompt во время retrieval и для последующей candidate generation, что уменьшает mismatch между construction dataset и evaluation. Главный недостаток состоит в том, что шаблон намеренно generic и поэтому не кодирует ни теорию юмора, ни информацию об аудитории.

\section{Инструкции для embedding queries}

Для адаптации embedding model к retrieval subtasks используются два коротких instruction templates:
\begin{quote}
    \ttfamily
    Instruct: Given a keyword for a joke, retrieve jokes would match the keyword.\\
    Query: \{keyword\}
\end{quote}

\begin{quote}
    \ttfamily
    Instruct: Given a prompt for a joke, retrieve jokes that would answer the prompt.\\
    Query: \{prompt\}
\end{quote}

Эти шаблоны важны, потому что pilot не использует raw keyword strings как embeddings напрямую. Он кодирует их внутри retrieval-oriented instructions, что делает semantic intent vector space более task-specific.

\section{Evaluation prompt}

System prompt для pairwise evaluation просит judge сравнить две шутки, написанные для одного prompt, и выбрать ровно одного winner. Инструкция явно не поощряет вознаграждать length или polish, если они не улучшают юмор. Такой дизайн намеренно прост и должен рассматриваться как baseline evaluation protocol, а не как финальная human-aligned rubric.

\chapter{Дополнительные замечания о программной архитектуре}

Репозиторий организован как набор asynchronous pipelines:
\begin{itemize}
    \item \texttt{JokesPipeline}: скачивает и предобрабатывает raw corpora;
    \item \texttt{EmbeddingsPipeline}: вычисляет dense joke embeddings;
    \item \texttt{KeywordsPipeline}: извлекает salient keywords через similarity и MMR;
    \item \texttt{ReferencesPipeline}: строит prompt-reference datasets через vector retrieval;
    \item \texttt{CandidatesPipeline}: генерирует шутки по prompts;
    \item \texttt{EvaluationPipeline}: выполняет pairwise judging и leaderboard aggregation.
\end{itemize}

Эта архитектура важна для финального диплома, потому что она отделяет data construction от model comparison. В результате последующий MRVF training stage можно добавить без переписывания остальной инфраструктуры. Отсутствующий компонент --- именно тот, который теперь становится основной целью future work: training pipeline, который потребляет построенный prompt-reference dataset и оптимизирует формальную objective, введенную в Chapter 2.
